\scanpage{134}
% lines 1–5
\[
\underline f^{m-n}: K_{*}(\underline f^{m})\qquad K_{*}(\underline f^{n})
\]
ou $m$ et $n$ sont des entiers rationnels, $m\quad n\quad 0$, d'ou des \deleted{\unsure{en|un}}homomorphismes, notes par le meme symbole:
\[
\begin{aligned}
& K^{*}(\underline f^{n},M)\quad K^{*}(\underline f^{m},M)\\
\text{et}\qquad & H^{*}(\underline f^{n},M)\quad H^{*}(\underline f^{m},M)\qquad (m\quad n\quad 0).
\end{aligned}
\]
% lines 6–13
Comme ces homomorphismes sont transitifs, on peut passer a la limite inductive dans ces homomorphismes, et poser
\[
\begin{aligned}
C^{*}((\underline f^{\deleted{\unsure{n}}}),M)&=\lim_{n} K^{*}(\underline f^{n},M)\\
\underline H^{*}((\underline f^{\deleted{\unsure{n}}}),M)\deleted{\unsure{*}}&=H^{*}(C^{*}((\underline f),M)=\lim_{n}H^{*}(\underline f^{n},M).
\end{aligned}
\]
(N.B. dans ces notations, $(\underline f)$ designe en fait l'ideal engendre dans $A$ par les $f_i$; on verra en effet ulterieurement que les objets introduits, \deleted{\unsure{ne dependent}} du moins $H^{*}((\underline f),M)$, ne depend aue de cet ideal, et meme seulement de la topologie $\underline m$-adique \deleted{\unsure{a}}qu'il definit, et non des $f_i$ eux-memes\unsure{.} 

% lines 14–22
Bien entendu, $C^{*}((\underline f),M)$ est un foncteur \deleted{exact} ($A$-lineaire exact) en $M$ et $H^{*}((\underline f),M)$ est un foncteur cohomologique en $M$.

Soit $\underline f$ comme ci-dessus, \deleted{\unsure{soit}} considere comme element du dual de $A^{n}$, soit $\underline g=(g_i)_{1\quad i\quad n}$ un element de $A^{n}$ lui-meme, on sait qu'on a alor\deleted{\unsure{x}}s, en designant par $e_{\underline g}$ l'operation de multiplication exterieure par $\underline g$ dans $A^{n}$, la \uline{formule} \uline{d'homotopie}
\[
i_{\underline f}e_{\underline g}+\unsure{e}_{\underline g}i_{\underline f}=\langle\underline g,\underline f\rangle\mathrm{id}.
\]
dans $A^{n}$, qui peut aussi s'ecrire, dans le complexe $K_{*}(\underline f)$, comme
\[
d\unsure{e}_{\underline g}+e_{\underline g}d=\underline g,\underline f\ \mathrm{id}.
\]
% lines 23–28
On en conclut la

\uline{Proposition 1}. Supposons que l'ideal $(\underline f)$ engendre par les $f_i$ soit $A$. \unsure{Alors} le complexe $K_{*}(\underline f)$ est homotopiquement equivalent a zero, et il en est \unsure{par} suite dememe des complexes $\unsure{K|H}_{*}(\underline f,M)$ et $K^{*}(\underline f,M)$ quels que soit $M$.
% slanted
\margin{\handwritten{\unsure{Autre|Plus|Ajouter un}, corollaire: les $f_i$ annulent les $H_{*}(\underline f,M)$, $H^{\cdot}(f,M)$}}

\uline{Corollaire} S'il en est ainsi, on a $H^{*}(\underline f,M)=0$ et $H^{*}((\underline f),M)=0$ pour \unsure{tout} $A$-module $M$.

% lines 29–38
Dans la description de $K^{*}(\underline f,M)$, il faut aussi ajouter \unsure{l'interpretation} des elements de degre $p$ de ce complexe comme des fonctions alternees de $p$ indices variables entre $1$ et $n$; la formule du cobord est alors formellement celle d'un cobord en cohomologie de Cech.

2. \uline{Cohomologie de Cech d'un recouvrement ouvert}.

Soient $X$ un preschema, $A=\mathrm{I}\ (X,\underline{\mathrm O}_{X})$, $\underline f=(f_i)_{1\quad i\quad n}$ un systeme de $n$ \unsure{elements} de $A$, $U_i=X_{f_i}$, $U=\quad U_i$, $\underline U$ le recouvrement de $U$ par les $U_i$. Soit $F$ un faisceau quasi-coherent sur $X$, posons
\[
M=\mathrm{I}\ (X,F)
\]
(N.B. les notations precedentes devraient etre deployees de meme !)
% archive sheet number: 133

